\documentclass[10pt,a4paper]{amsart}
\usepackage[margin=19mm]{geometry}
\usepackage{amsmath,amssymb,mathtools}
\usepackage[scr=rsfs,cal=euler]{mathalfa}
\usepackage{xfrac}
\usepackage[unicode,hidelinks,bookmarks=false]{hyperref}
\newcommand{\ie}{i.e.\ }
\newcommand{\cf}{cf.\ }
\newcommand{\resp}{respectively\ }
\newcommand{\Subsection}{\subsection}
\providecommand{\nobreakdash}{\nobreak\hbox{-}}
\setlength{\emergencystretch}{2em}
\multlinegap0pt
\usepackage{mathtools}
\usepackage{amssymb}
\usepackage[shortlabels]{enumitem}
\newtheorem{assumption}{Assumption}
\newtheorem{lemma}{Lemma}
\newcommand{\Leb}{\mathrm{Leb}}
\newcommand{\OT}{\mathrm{OT}}
\newcommand{\R}{\mathbb{R}}
\DeclareMathOperator{\gr}{gr}
\DeclareMathOperator{\spt}{spt}
\title{Quadratically Regularized Optimal Transport: Localization Bounds and Affine Case Analysis}
\author{Long Nguyen-Chi, Nam Nguyen, Binh Nguyen}
\date{Source: arXiv:2605.24644v2 (CC BY 4.0)}
\begin{document}
\maketitle

\noindent\emph{Source and licence.} Quadratically Regularized Optimal Transport: Localization Bounds and Affine Case Analysis. Version arXiv:2605.24644v2 (CC BY 4.0), distributed by arXiv under the Creative Commons Attribution 4.0 International licence, http://creativecommons.org/licenses/by/4.0/, as recorded in the arXiv licence metadata for this exact version and shown on the abstract page https://arxiv.org/abs/2605.24644v2. Full licence notice: \texttt{licenses/arxiv-2605.24644-CC-BY-4.0.txt}.\par\smallskip

\noindent\textit{Editorial scope.} Counted unit: the article's lemma lem:mse-lower-bound (source0.tex lines 812--822). The article's complete original proof of it is delivered with it (source0.tex lines 824--903, one \texttt{proof} environment of 80 lines). No mathematics is written, repaired or re-derived: the statement and the proof are the article's own lines, in the article's own notation, and no retained line is substituted or re-flowed.  Retained uncounted as the source-local context the proof needs: lines 269--274; lines 276--278.  Nothing was omitted from any retained span, so the package carries no omission record.  No retained line contains a citation, so the wrapper carries no bibliography and no bibliography entry.  No retained line contains a figure environment, an image-inclusion command or drawing code, so no image is delivered and the package declares no asset dependency.  Editorial formatting only, in addition: an AMS \texttt{amsart} preamble that restates the article's own declarations and macros used by the retained lines, namely the environments \texttt{assumption}, \texttt{lemma} on separate counters, together with the standard packages those lines need.  The retained environments print in this document as Assumption 1, Lemma 1 in order of appearance, whereas the article prints the same items with the numbering of its own full document.  The complete native source of the article as distributed in the arXiv e-print for this exact version is bundled unchanged as \texttt{sources/201291/source0.tex}.
\begin{assumption}\label{assu:mu}
There exists a bounded connected open Lipschitz domain $\Omega\subset\R^d$ with $\Omega\subset B(0,1)$ such that $\mu(dx)=\rho_\mu(x)\mathbf 1_{\Omega}(x)dx$, and there are constants $0<\underline\lambda_\mu\le\overline\lambda_\mu<\infty$ such that
$$
\underline\lambda_\mu\le\rho_\mu(x)\le\overline\lambda_\mu\text{ for Lebesgue-a.e. }x\in\Omega.
$$
\end{assumption}

\begin{assumption}\label{assu:nu}
The measure $\nu\in\mathcal{P}(\R^d)$ satisfies $\nu\ll\Leb$, $\spt\nu\subseteq B(0,1)$ and there exist constants $0<\overline\lambda_\nu<\infty$ such that $d\nu/dy\le\overline\lambda_\nu$ on $\spt\nu$.
\end{assumption}

\begin{lemma}[Optimality of the mean-squared bias scale]\label{lem:mse-lower-bound}
Assume Assumptions~\ref{assu:mu} and \ref{assu:nu} hold.
Let
\[
m_\varepsilon\coloneqq\int\|y-T(x)\|^2\pi_\varepsilon(dx,dy).
\]
Then, for all $\varepsilon\in(0,1]$, $m_\varepsilon\ge c_{\mathrm{mse}}\varepsilon^{\frac{2}{d+2}}$, where
\[
c_{\mathrm{mse}}\coloneqq\frac12\left(\frac{1}{96\overline\lambda_\nu\omega_d C_{\mathrm{val}}}\right)^{2/d}.
\]
\end{lemma}

\begin{proof}[Proof of Lemma~\ref{lem:mse-lower-bound}]
Let $h=d\pi_\varepsilon/d(\mu\otimes\nu)$ and set $Z(x,y)\coloneqq\|y-T(x)\|^2$.
We first prove that $m_\varepsilon>0$. 
The map $T$ is Borel, so $\gr T$ is a Borel subset of $\R^d\times\R^d$. 
If $m_\varepsilon=0$, then the nonnegative random variable $Z$ vanishes $\pi_\varepsilon$-a.s., and therefore $\pi_\varepsilon(\gr T)=1$.
On the other hand, since $\nu\ll\Leb$, every singleton has zero $\nu$-mass.
Hence
\[
(\mu\otimes\nu)(\gr T)=\int\nu(\{T(x)\})\mu(dx)=0.
\]
This contradicts $\pi_\varepsilon\ll\mu\otimes\nu$. 
Therefore $m_\varepsilon>0$.

Let $t=\sqrt{2m_\varepsilon}$. 
By Markov's inequality applied to $Z$,
\[
\pi_\varepsilon\bigl(\{(x,y):\|y-T(x)\|\le t\}\bigr)=\pi_\varepsilon(\{Z\le 2m_\varepsilon\})\ge\frac12.
\]
For $\nu$-a.e. $y$, define
$$
a(y)\coloneqq
\int_{T^{-1}(B(y,t))}h(x,y)\mu(dx).
$$
Then $0\le a(y)\le 1$ for $\nu$-a.e. $y$, and
$$
\int a(y)\nu(dy)
=
\pi_\varepsilon(\{(x,y):\|y-T(x)\|\le t\})
\ge \frac12.
$$
Let $D=\{y:a(y)\ge 1/4\}$. 
Since $a\le 1$,
$$
\frac12
\le
\int a(y)\nu(dy)
\le
\nu(D)+\frac14(1-\nu(D)),
$$
and therefore $\nu(D)\ge 1/3$.

For every $y\in D$, Cauchy-Schwarz yields
$$
\frac{1}{16}
\le
\left(\int_{T^{-1}(B(y,t))}h(x,y)\mu(dx)\right)^2
\le
\mu(T^{-1}(B(y,t)))
\int_{T^{-1}(B(y,t))}h(x,y)^2\mu(dx).
$$
Since $T_\#\mu=\nu$ and Assumption~\ref{assu:nu} gives $\nu(B(y,t))\le\overline\lambda_\nu\omega_dt^d$, we obtain
\[
\int h(x,y)^2\mu(dx)\ge\frac{1}{16\overline\lambda_\nu\omega_dt^d}\text{ for every }y\in D.
\]
Integrating over $D$ gives
$$
\|h\|_{L^2(\mu\otimes\nu)}^2
\ge
\frac{1}{48\overline\lambda_\nu\omega_dt^d}.
$$
Since
$$
\Delta_\varepsilon
=
\int cd\pi_\varepsilon-\OT(\mu,\nu)
+
\frac{\varepsilon}{2}\|h\|_{L^2(\mu\otimes\nu)}^2
\ge
\frac{\varepsilon}{2}\|h\|_{L^2(\mu\otimes\nu)}^2,
$$
we obtain
$$
\Delta_\varepsilon\ge\frac{\varepsilon}{96\overline\lambda_\nu\omega_dt^d}=\frac{\varepsilon}{96\overline\lambda_\nu\omega_d(2m_\varepsilon)^{d/2}}.
$$
Combining this lower bound with $\Delta_\varepsilon\le C_{\mathrm{val}}\varepsilon^{2/(d+2)}$ gives
$$
m_\varepsilon^{d/2}\ge\frac{1}{96\times2^{d/2}\overline\lambda_\nu\omega_d C_{\mathrm{val}}}\varepsilon^{d/(d+2)}.
$$
Raising both sides to the power $2/d$ gives the claimed constant $c_{\mathrm{mse}}$.
\end{proof}

\end{document}
